\chapter{\MakeUppercase{Methodology and Testing}}

\section{Experimental Methodology}
The evaluation compares Epidemic, Spray and Wait, and Spatio-Semantic routing under three traffic levels. Each protocol is executed for 20 runs per traffic level. The simulator averages the resulting metrics and writes one \ac{csv} file per protocol. This produces a direct comparison because all protocols are evaluated using the same node counts, area size, buffer size, contact range, message \ac{ttl}, and traffic probabilities.

\section{Traffic Scenarios}
The traffic scenarios are defined by per-node message generation probability:

\begin{table}[h!]
	\centering
	\caption{Traffic scenarios}
	\renewcommand{\arraystretch}{1.3}
	\begin{tabular}{|l|l|l|}
		\hline
		\textbf{Traffic Level} & \textbf{Generation Probability} & \textbf{Interpretation} \\
		\hline
		Low & 3/3600 & Light background communication \\
		\hline
		Medium & 8/3600 & Moderate disaster coordination traffic \\
		\hline
		High & 20/3600 & Congested emergency communication \\
		\hline
	\end{tabular}
	\label{tab:traffic_scenarios}
\end{table}

\section{Baseline Protocol Testing}
Epidemic routing is tested to measure the behavior of unrestricted replication. It forwards every message to every encountered node that does not already have it, provided the receiver buffer has space. This baseline helps show the cost of blind reachability.

Spray and Wait routing is tested to measure copy-limited forwarding. The implementation gives one copy at a time during the spray phase and forwards a single-copy message only when the destination is encountered. This baseline helps show the cost of conservative forwarding.

\section{Proposed Protocol Testing}
Spatio-Semantic routing is tested using the same traffic and mobility conditions. Its main test objective is to determine whether spatial, semantic, and priority-aware decisions improve performance under load. The protocol is expected to be especially strong when traffic is high because congestion exposes whether buffer and relay choices are meaningful.

\section{Dashboard and Visual Testing}
The project includes a Streamlit dashboard for plotting metric trends, radar comparisons, raw results, and launching the live visual simulation. The live visualizer compares the three protocols side by side with the same initial seed and displays delivery, critical delivery, critical delay, overhead, drops, contacts, node roles, and critical-message carriers.

\begin{figure}[h!]
    \centering
    \fbox{\parbox{0.88\textwidth}{\centering
    Image placeholder: add a screenshot of the Streamlit Metric Trends tab showing delivery ratio, critical delivery ratio, delay, overhead, and buffer drops.}}
    \caption{Dashboard metric trends for the three routing protocols}
    \label{fig:dashboard_trends_placeholder}
\end{figure}

\begin{figure}[h!]
    \centering
    \fbox{\parbox{0.88\textwidth}{\centering
    Image placeholder: add a screenshot of the Performance Radar tab for the high traffic scenario.}}
    \caption{Radar comparison of multi-objective routing performance}
    \label{fig:radar_placeholder}
\end{figure}

\begin{figure}[h!]
    \centering
    \fbox{\parbox{0.88\textwidth}{\centering
    Image placeholder: add a screenshot of the three-panel live simulation with Epidemic, Spray and Wait, and Spatio-Semantic views.}}
    \caption{Live three-panel simulation of competing DTN routing protocols}
    \label{fig:live_sim_placeholder}
\end{figure}

\section{Testing Criteria}
The protocol is considered successful if it satisfies the following criteria:

\begin{itemize}
	\item It should match or exceed Epidemic routing in delivery under light and medium traffic without increasing congestion unnecessarily.
	\item It should significantly outperform Epidemic routing under high traffic, where flooding is expected to collapse.
	\item It should reduce critical-message delay compared with Spray and Wait, because emergency messages must not wait passively for direct delivery.
	\item It should maintain a clear advantage in critical delivery ratio under congested conditions.
\end{itemize}
